% Manuscript-ready draft. Values are copied from frozen CARE 2.0 trace and
% summary artifacts; no values in this section are illustrative.

\subsection{How source evidence changes sequential decisions}

Aggregate performance alone does not show how transferred knowledge altered an
experimental trajectory. We therefore inspected three representative routes: a
successful molecular-property override, a successful materials consensus
decision, and a negative-transfer case. In each example, all unqueried target
outcomes remained hidden from the language model and the controller until the
corresponding candidate was executed.

\paragraph{A source-grounded override produced an early molecular-property gain.}
For FreeSolv $\rightarrow$ Lipophilicity, the three initial target observations
gave a best value of 67.25. In the first online round, the target-only Gaussian
process ranked \texttt{lipo\_2256} first. The proposer and critic instead selected
\texttt{lipo\_0635}, which was ranked fourth by the GP and had an expected
improvement of 3.90. The recorded rationale connected the candidate's
halogen-rich, low-polarity structure to the target observations rather than
copying the source outcome directly. The candidate was then executed and
revealed a target value of 86.875, increasing the best observed value by 19.625
in one round. Across the ten-round trajectory, the best-so-far AUC was 89.575,
compared with 83.250 for the same-initial target-only GP comparator
($\Delta$AUC=+6.325). This example shows the intended role of the language
model: it can overrule the default acquisition ranking when source evidence and
target observations support a specific, testable alternative.

\paragraph{Independent target and source signals converged on a high-value material.}
For experimental band gap $\rightarrow$ dielectric response, the initial target
best was 24.2902. In round zero, the five GP-eligible candidates had very similar
posterior means and expected improvements. The proposer and critic chose the
tellurium candidate \texttt{real\_matbench\_dielectric\_04424}. It was only fifth
in the target GP ranking, but first in the combined consensus ranking because
the source-derived polarizability hypothesis and the target posterior both
supported it. Its target value, hidden at decision time, was 59.5114. The online
trajectory consequently attained a best-so-far AUC of 59.5114, compared with
52.7663 for the same-initial target-only GP ($\Delta$AUC=+6.7452). The final best
was identical between the two methods; the gain came from finding the same
high-value region earlier, which is precisely what the AUC reward measures.

\paragraph{Persistent transfer can also delay discovery.}
The experimental band gap $\rightarrow$ Materials Project band-gap route is a
negative control. Its initial best was already 75.5362. None of the ten online
LLM selections improved that value, whereas the same-initial target-only GP
eventually reached 100.0. The corresponding AUC difference was $-2.4464$ and
the final-best difference was $-24.4638$. The trace shows that the controller
continued source transfer in every round, even after repeated lanthanide-halide
candidates failed to improve the target best. This failure is informative: a
plausible source-grounded mechanism is not sufficient evidence for continued
transfer. The result motivates a stricter sequential stopping rule based on
accumulated target contradictions rather than on the coherence of the initial
hypothesis alone.

Together, these cases support a narrower claim than universal cross-domain
transfer. CARE 2.0 can convert source evidence into actionable deviations from
a target-only acquisition policy, and those deviations can substantially
improve early discovery on some routes. The same mechanism can also preserve a
misleading prior for too long. Reporting both outcomes makes the calibration
problem explicit and separates the contribution of hypothesis generation from
the still-unresolved problem of reliable transfer rejection.

% Frozen evidence sources:
% experiments/care_replay/results/2026-08-16-opus48-bounded-ei-development-v2/
%   molecular_freesolv_to_lipophilicity/{summary.json,llm_trace.jsonl}
%   materials_expt_gap_to_dielectric/{summary.json,llm_trace.jsonl}
%   materials_expt_gap_to_mp_gap/{summary.json,llm_trace.jsonl}

